\begin{satz}{Existenz eingeschränkter Fast--Isometrien}
            {ExistenzEingeschraenkterFastIsometrien}
Seien $m, N \in \N$,
$((\Omega, \mfa, \Prim), (\gamma, \delta, \widetilde g), (N, m, mN))$
ein unabhängiges $(N(0,1),3)$--Tupel (zu $(3, (N, m, mN))$) und
$g\coloneqq \Mat_{m,N} \circ \widetilde g$.\\
Seien $(E, \norm\cdot_E)$ ein normierter $\R$--Vektorraum,
$0\neq x \in E^n$ und $\emptyset \neq T \subseteq S^{m-1}$.\\
Seien $G_x$ der Gauß'sche Operator zum Tupel $(x,g)$ 
(vgl.~\ref{GaussscherOperator}) und $(Z_{\gamma,x}, M_{\delta, T})$ das
Auswertungspaar zu $(\gamma, \delta, x, T)$ (vgl.~\ref{Auswertungspaar}).\\
Es gelte:
\begin{enumerate}[(a)]
	\item $\norm\cdot_E \circ \bprod \cdot x _N \le \norm\cdot_{2, N}$,
	\item $T = -T$,
	\item $\E\left(M_{\delta, T}\right) < \E\left(Z_{\gamma, x} \right)$.
\end{enumerate}
Dann gilt $\E(Z_{\gamma, x}) > 0$, $\E(M_{\delta, T}) \ge 0$ und es gibt 
$\mca \subseteq L(\ell_2^m, (E, \norm\cdot_E))$ mit:
\begin{enumerate}[(1)]
  \item\label{ExistenzEingeschraenkterFastIsometrien1}
  $\mca \subseteq \Bild(G_x)$, es gibt $Q \in \mfa$ mit $\Prim(Q)>0$ und
  $Q \subseteq G_x\inv(\mca)$, insbesondere also $\mca \neq \emptyset$.
	\item \label{ExistenzEingeschraenkterFastIsometrien2}
	$\forall\, A \in \mca: \inf\menge{\norm{At}_E}{t \in T} > 0$\\
	Falls $T = S^{m-1}$, so ist jedes $A \in \mca$ injektiv.
	\item \label{ExistenzEingeschraenkterFastIsometrien3}
	$\forall\, A \in \mca:
	\dfrac{\sup\menge{\norm{At}_E}{t \in T}}{\inf\menge{\norm{At}_E}{t \in T}}
	\le \dfrac{\E\left(Z_{\gamma, x} \right) + \E\left(M_{\delta, T}\right)}
	{\E\left(Z_{\gamma, x} \right) - \E\left(M_{\delta, T}\right)}
	=\dfrac{1 + \frac{\E(M_{\delta, T})}{\E(Z_{\gamma, x})}}
	       {1 - \frac{\E(M_{\delta, T})}{\E(Z_{\gamma, x})}}$\\
	Ist $T = S^{m-1}$, so gilt für jedes $A \in \mca$: 
	$A\inv\in \Inv(A(\R^m),\R^m)$ und
	\[\opnorm A {\R^m} E \opnorm{A\inv}{A(\R^m)}{\R^m}
	\le \dfrac{1 + \frac{\E(M_{\delta, T})}{\E(Z_{\gamma, x})}}
	       {1 - \frac{\E(M_{\delta, T})}{\E(Z_{\gamma, x})}}\text.\]
\end{enumerate}
Ist zusätzlich $\eps \in\!\  ]0, 1[$ mit
$\E\left(M_{\delta, T}\right) \le \eps\E\left(Z_{\gamma, x} \right)$, so gilt
auch
\begin{enumerate}
\item[(3')] \label{ExistenzEingeschraenkterFastIsometrien3strich}
$\forall\, A \in \mca:
	\dfrac{\sup\menge{\norm{At}_E}{t \in T}}{\inf\menge{\norm{At}_E}{t \in T}}
	\le \dfrac{1 + \eps}	{1 - \eps}$\\
	Gilt $T = S^{m-1}$, so ist für jedes $A \in \mca$:
	\[\opnorm A {\R^m} E \opnorm{A\inv}{A(\R^m)}{\R^m}
	\le \dfrac{1 + \eps}	{1 - \eps}\text.\]
\end{enumerate}
\end{satz}
\begin{proof}
Zunächst sind offenbar die Voraussetzungen von~\ref{GordonschePaareUndSubeuklidischeBildhalbnormen} erfüllt.\\
Seien
\[f\coloneqq  \inf_{t \in T} \, \norm \cdot_E \circ ev(t, \cdot) \circ G_x
\quad \text{ und }\quad
  h\coloneqq  \sup_{t \in T} \, \norm \cdot_E \circ ev(t, \cdot) \circ G_x.\]
Nach~\refclaim{IntegrierbarkeitVonExtremaInGordonschenPaaren}
{IntegrierbarkeitVonExtremaInGordonschenPaaren2} sind $f, h$ integrierbar.\\
Wegen (b) und~\refclaim{GordonschePaareUndSubeuklidischeBildhalbnormen}
{GordonschePaareUndSubeuklidischeBildhalbnormen2} gilt weiter:
\begin{align} \E\left(\inf_{t\in T}\bprod t \delta_m \right)
	= - \E\left(\sup_{t\in T}\bprod t \delta_m \right) = 
	- \E\left(M_{\delta,T}\right)\text.\tag {$\ast$}
\end{align}
Somit gilt wegen der Integrierbarkeit von $Z_{\gamma, x}$ und $M_{\delta, T}$
(siehe~\ref{Auswertungspaar}):
\begin{align*}
0 & \overset{(c)}< \E\left(Z_{\gamma, x}\right) - \E\left(M_{\delta, T} \right)
\overset{(\ast)}= \E\left(Z_{\gamma, x}\right) 
+ \E\left(\inf_{t\in T}\bprod t \delta_m \right)\\
&\overset{\refclaim{GordonschePaareUndSubeuklidischeBildhalbnormen}
{GordonschePaareUndSubeuklidischeBildhalbnormen1}}\le \E(f) \le \E(h)
\overset{\refclaim{GordonschePaareUndSubeuklidischeBildhalbnormen}
{GordonschePaareUndSubeuklidischeBildhalbnormen1}}\le
\E\left(Z_{\gamma, x}\right) + \E\left(M_{\delta, T} \right)\text.
\end{align*}
Somit $\E(Z_{\gamma,x}) - \E(M_{\delta, T}) \le
\E(Z_{\gamma,x}) + \E(M_{\delta, T})$ und daher $\E(M_{\delta, T}) \ge 0$, also
gilt wegen $(c)$ auch $\E(Z_{\gamma,x})>\E(M_{\delta, T}) \ge 0$.\\
Seien nun $a\coloneqq \E(Z_{\gamma,x}) - \E(M_{\delta, T})$ und
$b\coloneqq \E(Z_{\gamma,x}) + \E(M_{\delta, T})$.\\
Wegen $f \ge 0$ und $h \ge 0$ und nach obiger Rechnung gilt:
\[ 0 < a \le \E(f) \le E(h) \le b\text.\]
Mit der Integrierbarkeit von $f$ und $h$ sind somit die Voraussetzungen von~\ref{Erwartungswertlemma} für $f, h, a, b$ erfüllt und für
$Q \coloneqq \menge{\omega \in \Omega}{f(\omega)> 0 \wedge ah(\omega) \le bf(\omega)}$
gilt daher (gemäß~\ref{Erwartungswertlemma}): $Q\in \mfa$ und $\Prim(Q) > 0$.\\
Sei nun $\mca\coloneqq G_x(Q)$.
Dann ist $G_x\inv(\mca) \supseteq Q$.\\
Sei $A \in \mca$. Dann existiert ein $\omega \in Q$ mit $A = G_x(\omega)$.\\
Wegen $\omega \in Q$ ist also
\[0 < f(\omega) = \inf_{t \in T}\norm{G_x(\omega)t}_E 
    = \inf_{t \in T}\norm{At}_E \]
und außerdem $ah(\omega) \le bf(\omega)$, was wegen $f(\omega)> 0$ und $a >0$
auch $\frac{h(\omega)}{f(\omega)} \le \frac b a$ bedeutet, sodass mit
$\E(Z_{\gamma, x})>0$ gilt:
\begin{align*}
&\,\frac{\sup\menge{\norm{At}_E}{t \in T}}
       {\inf\menge{\norm{At}_E}{t \in T}}
 =\frac{\sup\menge{\norm{G_x(\omega)t}_E}{t \in T}}
       {\inf\menge{\norm{G_x(\omega)t}_E}{t \in T}}
 = \frac{h(\omega)}{f(\omega)}\\ \le & \,\frac b a = 
 \frac{\E(Z_{\gamma,x})+\E(M_{\delta, T})}{\E(Z_{\gamma,x})-\E(M_{\delta, T})}
 = \frac{\E(Z_{\gamma,x})(1 + \frac{\E(M_{\delta, T})}{\E(Z_{\gamma,x})})}
        {\E(Z_{\gamma,x})(1 - \frac{\E(M_{\delta, T})}{\E(Z_{\gamma,x})})}
 = \frac{1 + \frac{\E(M_{\delta, T})}{\E(Z_{\gamma,x})}}
        {1 - \frac{\E(M_{\delta, T})}{\E(Z_{\gamma,x})}}\text.
 \end{align*}
Ist $T = S^{m-1}$, so folgt mit $\inf\menge{\norm{At}_E}{t \in S^{m-1}} > 0$
und $S^{m-1} = \kran{\ell^2_m}01$ nach~\ref{InjektivitätLinearerAbbildungenZwischenNormiertenRäumen} die Injektivität
von $A$.\\
Weiter ist der Raum 
$(A(\R^m), \restrict{\norm\cdot_E}{A(\R^m)})$ ein endlichdimensionaler
normierter Raum und damit ein Banachraum.\\
Fasse nun $A$ als Abbildung $A:\R^m \rightarrow A(\R^m)$ auf.\\
Dann ist $A$ bijektiv, linear und stetig, sodass nach~\ref{OperatornormDerInversen} gilt:
\[\opnorm{A\inv}{A(\R^m)}{\R^m} 
= \frac 1 {\inf\menge{\norm{At}_E}{t \in S^{m-1}}}\text.\]
Wegen $\opnorm A {\R^m} E = \sup\menge{\norm{At}_E}{t \in S^{m-1}}$ gilt also
\[\opnorm A {\R^m} E \opnorm{A\inv}{A(\R^m)}{\R^m}
= \frac{\sup\menge{\norm{At}_E}{t \in S^{m-1}}}
       {\inf\menge{\norm{At}_E}{t \in S^{m-1}}}\text.\]
Damit sind $(1)$, $(2)$, $(3)$ gezeigt.\\
Sei nun $\eps\in\!\ ]0, 1[$ mit $\E(M_{\delta, T})\le \eps \E(Z_{\gamma,x})$.\\
Wegen $\E(Z_{\gamma,x}) > 0$ ist also 
$\frac{\E(M_{\delta, T})}{\E(Z_{\gamma,x})} \le \eps$ und daher
\[\frac{1 + \frac{\E(M_{\delta, T})}{\E(Z_{\gamma,x})}}
       {1 - \frac{\E(M_{\delta, T})}{\E(Z_{\gamma,x})}}
\le \frac{1 + \eps}{1 - \frac{\E(M_{\delta, T})}{\E(Z_{\gamma,x})}}
\le \frac{1 + \eps}{1 - \eps} \text.\]
Wegen $(3)$ ist damit auch $(3')$ gezeigt.
\end{proof}